%*****************************************************************************
% Layout-Vorschlag zur Erstellung von Bachelor-, Master oder Diplomarbeiten  *
% Modifikation moeglich                                                      *
% S. Gippert, 2012                                                           *
% ****************************************************************************

\documentclass[11pt,a4paper]{scrbook}
\usepackage{graphicx}
\usepackage{ngerman}
\usepackage[ngerman]{babel}
\usepackage{amsmath}
\usepackage{amssymb}
\usepackage{amsthm}


\usepackage{multirow}
\usepackage{booktabs}

\usepackage{hyperref}

% Umgebungen
\renewcommand{\baselinestretch}{1.50}\normalsize
\newenvironment{Beweis}{\begin{proof}[Beweis]}{\end{proof}}
\newtheorem{Satz}{Satz}[chapter]% Definitionen
\def\Real{{{\rm I\!R}}}

% Dokumentbeginn
\begin{document}

% Titelseite ----------ANPASSEN!----------------------
\begin{titlepage}
\begin{center}
\huge \textbf{{Layerwise Relevance Propagation For Deep Neural Networks}} \\
\vspace{2cm}
\LARGE\textbf{{Abschlussbericht}}\\
\vspace{2cm}
\normalsize


vorgelegt von  \textbf{Gruppe 4\\ Theo Conrads, Robin Kühling, Marc Bremser}\\
am \today \\
\vspace{2.5cm}
\textbf{am Mathematischen Institut der \\Universität zu Köln}\\
\vspace{3cm}
\end{center}
\normalsize{
\begin{center}
Erstgutachter:  {Prof.~Dr.~Axel Klawonn} \\
\vfill
\includegraphics[scale=0.25]{uzk_logo} % bunt
% \includegraphics[scale=0.25]{uzk_logo_grau} % grau
\end{center}
}
\end{titlepage}


\pagestyle{plain}
\cleardoublepage
\setcounter{page}{1}
\pagenumbering{roman}
\tableofcontents

\listoffigures
\addcontentsline{toc}{chapter}{Abbildungsverzeichnis}

\listoftables
\addcontentsline{toc}{chapter}{Tabellenverzeichnis}

\cleardoublepage

% Dateien
\chapter{Einleitung}
Hier grundsätzliche Einführung
\input{Kurzeinführung_CNN}
\chapter{Deep Taylor Decomposition}
\section{Taylor Decomposition}
Einführung in die einfache Taylor Decomposition
\section{Relevance Models}
\subsection{Das Min-Max-Model}
Theorie zum Min-Max-Model
\subsection{Das Training-Free Model}
Theorie zum Training-Free Ansatz
\chapter{Layerwise Relevance Propagation}
\section{Herleitung der Basisregel}
Herleitung nach PLOS-Paper. Bezug zu DTD herstellen
\section{Erweiterungen der Basisregel}
LRP-$\epsilon$, LRP-$\alpha-\beta$, LRP-$\gamma$
\chapter{Implementierung}
\section{Implementierung des Min-Max-Models}
Erklärung der wichtigsten Codeabschnitte.
Ergebnisse werden voraussichtlich erst im Kapitel Ergebnisse präsentiert.
\section{Implementierung LRP}
Erklärung der wichtigsten Code-Abschnitte
\chapter{Ergebnisse}
\section{Ergebnisse Min-Max-Model}
Präsentation der Ergebnisse. Bilder/Plots (Zeitaufwand etc.)

\section{Ergebnisse LRP inkl. Training-Free}
Präsentation der Ergebnisse. Bilder/Plots (Zeitaufwand etc.)
\chapter{Fazit}
Erklärung warum das ganze Thema großer Schrott ist :-)

% Literaturverzeichnis
\bibliography{thesisbib}{}
\bibliographystyle{plain}
\addcontentsline{toc}{chapter}{Literaturverzeichnis}
%%\newpage
\thispagestyle{empty}
\chapter*{Erkl\"arung}
\addcontentsline{toc}{chapter}{Erkl\"arung}
\noindent
Hiermit versichere ich an Eides statt, dass ich die vorliegende Arbeit selbstst\"andig und ohne die Benutzung anderer als der angegebenen Hilfsmittel angefertigt habe. Alle Stellen, die w\"ortlich oder sinngem\"a\ss   aus ver\"offentlichten und nicht ver\"offentlichten Schriften entnommen wurden, sind als solche kenntlich gemacht. Die Arbeit ist in gleicher oder \"ahnlicher Form oder auszugsweise im Rahmen einer anderen Pr\"ufung noch nicht vorgelegt worden. Ich versichere, dass die eingereichte elektronische Fassung der eingereichten Druckfassung vollst\"andig entspricht.
\begin{tabular}{p{5cm} p{7cm}}
\vspace{4cm}\\
\hline
Ort, Datum & Unterschrift
\end{tabular}


\end{document}
